Tato krátká sekce shrnuje kroky pro omezení sběru a bezpečný datový tok; uplatněte minimalizaci, účelové omezení a anonymizaci/pseudonymizaci \cite{kb_ai_101_tipu_pdf_04e4a115_page_8_1}.

\begin{itemize}
  \item Definujte účel polí; sbírejte jen nezbytné údaje, používejte pseudonymy/ID místo jména/fotek \cite{kb_ai_101_tipu_pdf_04e4a115_page_8_1}.
  \item Anonymizujte/pseudonymizujte před sdílením a omezte přístupy; zabraňte nahrávání citlivých souborů do neověřených cloudů, používejte interní nebo enterprise úložiště \cite{kb_ai_101_tipu_pdf_04e4a115_page_8_1,kb_ai_101_tipu_pdf_04e4a115_page_53_2}.
  \item U pilotů prověřte DPIA a DPA s poskytovatelem; stanovte retenci a role human-in-the-loop \cite{kb_malinka_chatgpt_ve_vyuce_dva_roky_pote_2024_pdf_90bc3aaf_page_12_1}.
\end{itemize}

Krátký doporučený datový tok:
\begin{enumerate}
  \item Sběr (minimální pole) → předzpracování (kontrola PII, anonymizace) → šifrované interní/enterprise uložení.
  \item Zpracování: AI pracuje jen s anonymizovanými daty v rámci účelu; retence a trvalé mazání záloh podle plánu.
\end{enumerate}

Praktické příklady: formativní kvíz — jen anonymní skóre a čas; programátorský projekt — školní repozitář; při externích AI službách vyžádejte souhlas a odstranění PII \cite{kb_ai_101_tipu_pdf_04e4a115_page_53_2}.